\section{国际民事诉讼}
\begin{enumerate}
  \item 管辖权。民诉通用。涉外纠纷我们是扩大管辖的。6 类。
        \begin{enumerate}
          \item 牵连管辖。合同、财产权益纠纷 + 被告在中国没有住所 = \zd{沾边就管}。
          \item 书面协议管辖法院。\zd{实际联系原则}。只要和案件有一点实际联系。特别爱考。
                \begin{enumerate}
                  \item 仅限合同、财产权益纠纷。
                  \item 选法院。
                        \begin{enumerate}
                          \item 受\zd{实际联系原则}限制。\al{最密切联系原则 \cuo}。
                          \item 不得违反级别和专属管辖的规定。\al{不动产}。
                          \item \zd{海诉}就近便利原则，不受实际联系原则限制。
                        \end{enumerate}
                  \item \zd{选（实体）法律}不受实际联系原则限制；一审法庭辩论终结前。\al{约定适用英国法 \dui，约定适用中国民事诉讼法 \cuo}。
                \end{enumerate}
          \item 不方便法院/管辖原则。
                \begin{enumerate}
                  \item 涉外\zd{同时}满足 6 个条件，\zd{可以}驳回原告起诉，告知向更方便的外国法院起诉。
                  \item 考试数有没有 6 个。\al{波兰甲公司认为本案由波兰法院管辖更为方便，可裁定驳回起诉 \cuo}。
                \end{enumerate}
          \item 专属管辖。民诉法第 273 条。\zd{在中国履行的}中外合资、中外合作、中外合作勘探开发自然资源，\zd{只能适用中国法}，但不排除可以用仲裁方式解决。
          \item 平行诉讼和一事再诉。国际私法特有，考试重复率高。
                \begin{enumerate}
                  \item \zd{不排斥}同一涉外民商事纠纷在不同国家的\zd{平行诉讼。平行审平行判}。
                  \item \zd{在先原则}（避免冲突判决的产生）。我们\zd{已判}外国判决拿来不认，外国判决我们\zd{已认}，一事再诉不理。
                        \al{Y 国人甲和我国乙在 Y 诉讼离婚，我国法院受理了 Y 执行判决的申请，乙向同一法院起诉离婚}。
                \end{enumerate}
          \item 国际商事法庭。一带一路背景。
                \begin{enumerate}
                  \item 级别。是\zd{最高法}的常设审判机构。一审终审，不能上诉，满足条件\zd{可以申请再审}。
                  \item 3 大类。
                        \begin{enumerate}
                          \item 当事人书面协议选它，且标的额为\zd{人民币 3 亿元以上}的\zd{第一审国际商事案件}。
                          \item 申请保全、撤销、执行国际商事仲裁裁决。
                          \item 全国范围内有重大影响 \yjt 高院搞不定报请最高法 \yjt 最高法认为应当管的。
                        \end{enumerate}
                  \item 证据的特殊规定。一强调两突破。
                        \begin{enumerate}
                          \item 强调。域外证据均\zd{应当}在法庭上\zd{质证}。
                          \item \zd{仅限英文}证据 + 对方同意 = 可以不提交翻译件。
                          \item 取证、质证可以采用信息网络方式实现。
                        \end{enumerate}
                  \item 判决。少数服从多数，\zd{少数意见可以在裁判文书中载明}。
                \end{enumerate}
          \item \zd{关于涉外民商事案件管辖若干问题的规定}。新增。
                \begin{enumerate}
                  \item \zd{下沉管辖为原则，集中管辖为例外}。
                  \item 标的额管辖标准。
                        \begin{enumerate}
                          \item 高院 50 亿以上。
                          \item 中院。有钱的 4000W+。4 大直辖市 + 东部沿海 5 省（山东、江苏、浙江、福建、广州）。没钱的 2000W+。其他省和自治区。
                        \end{enumerate}
                  \item 高院报批最高法后可指定基层 \hz 中院实施区域集中管辖。考不着。
                  \item 例外。涉外\zd{海事}海商、涉外\zd{知}识产权、涉外\zd{生态环境民事赔偿}、涉外\zd{环境公益诉讼}。
                \end{enumerate}
        \end{enumerate}
  \item 在中国无住所的（\zd{只看住所不看国籍}）。特殊的上诉期、答辩期、对上诉状的答辩期 30 天、特殊的送达方式。
  \item \zd{审期不受限制}、审理程序可以是普通程序 \hz 简易程序、\zd{不适用小额诉讼}。
  \item 外国人的民事诉讼地位。
        \begin{enumerate}
          \item 当事人的身份证明。
                \begin{enumerate}
                  \item 自然人。护照。
                  \item 法人。企业组织 + 代表参加诉讼的人。有/没有外交关系 \yjt 1 公证加 1/2 认证。
                \end{enumerate}
          \item 委托律师和诉讼代理人的限制。
                \begin{enumerate}
                  \item 委托律师只能委托\zd{我国律师以律师身份}参与诉讼。
                  \item \zd{可以委托本国律师}以非律师身份担任诉讼代理人。
                  \item \zd{使领馆官员}间接介入诉讼\zd{看身份}。外交代表/个人身份 \yjt 有/无特权豁免。
                        \begin{kk}{授权委托书的证明手续}
                          \begin{enumerate}
                            \item \zd{法官见证下签署即有效}。
                            \item 中国境内签署附公证书。
                            \item 中国境外签署\zd{要求同法人}。
                          \end{enumerate}
                        \end{kk}
                \end{enumerate}
        \end{enumerate}
\end{enumerate}
